\section{Method}

\subsection{Benchmark objective}

The benchmark is designed to compare mechanistic model families for weekly
influenza hospitalization-rate forecasting under a common data-processing,
fitting, and evaluation protocol. The objective is twofold:
\begin{enumerate}[leftmargin=1.5em]
\item quantify point-forecast and uncertainty performance across manual and
discovered epidemic models; and
\item determine whether constrained non-LLM structure discovery provides
useful age-specific latent or observation structures beyond strong hand-designed
baselines.
\end{enumerate}

The benchmark treats all epidemic states in normalized population units and
keeps the total latent population approximately conserved through explicit
penalties. All models expose a single observation scaling coefficient $\rho$ so
that the benchmark compares latent disease dynamics and observation semantics
rather than arbitrary output rescaling tricks.

\subsection{Data processing and series construction}

The raw FluSurv-NET CSV is read with \texttt{skiprows=2} to discard metadata
lines above the table. Rows with missing \texttt{WEEK} or \texttt{YEAR.1} are
removed, the remaining rows are sorted chronologically by year and week, and a
continuous integer time index $t \in \{0, \dots, T-1\}$ is created. The primary
series is the \emph{Entire Network} / \emph{Overall} weekly rate. Robustness is
evaluated on five age-stratified series:
\texttt{0--4 yr}, \texttt{5--17 yr}, \texttt{18--49 yr},
\texttt{50--64 yr}, and \texttt{$\geq$65 yr}.

\subsection{Manual model families}

\paragraph{Deterministic SEIR.}
The deterministic baseline uses seasonal transmission
\[
\beta_t = \operatorname{softplus}\!\left(b_0 +
b_1 \sin(2 \pi t / 52) + b_2 \cos(2 \pi t / 52)\right)
\]
with discrete latent dynamics over $(S,E,I,R)$ and observation model
$\hat{y}_t = \rho I_t$. Trainable parameters include the seasonal coefficients,
transition rates, observation scale, and initial exposed/infectious mass.

\paragraph{Probabilistic SEIR.}
The probabilistic baseline shares the same latent dynamics but replaces the
deterministic observation equation with a Student-$t$ observation model. The
benchmark fits this model by MAP estimation and then constructs bootstrap-based
raw predictive intervals. Point forecasts are the bootstrap centers; interval
calibration is handled later as benchmark-level post-processing.

\paragraph{Fractional SEIR.}
The fractional baseline extends SEIR with a shared order parameter
$\alpha \in [0.7, 1.0]$ and a discrete memory approximation. The intent is to
test whether long-memory latent dynamics provide stable forecasting benefit for
the hospitalization-rate target.

\paragraph{Hospitalized SEIHR.}
The hospitalization-aware manual baseline uses five compartments
$(S,E,I,H,R)$ with both direct recovery and hospitalization pathways:
\[
S \rightarrow E \rightarrow I,\qquad
I \rightarrow R,\qquad
I \rightarrow H \rightarrow R.
\]
Its weekly discrete dynamics are
\begin{align*}
S_{t+1} &= S_t - \beta_t S_t I_t, \\
E_{t+1} &= E_t + \beta_t S_t I_t - \sigma E_t, \\
I_{t+1} &= I_t + \sigma E_t - \eta I_t - \gamma_i I_t, \\
H_{t+1} &= H_t + \eta I_t - \gamma_h H_t, \\
R_{t+1} &= R_t + \gamma_i I_t + \gamma_h H_t,
\end{align*}
with observation model $\hat{y}_t = \rho H_t$. This corrected specification
avoids forcing every infectious individual through the hospitalized compartment.

\paragraph{Delayed-observation SEIR.}
This baseline preserves the standard SEIR latent dynamics but changes the
observation equation to
\[
\hat{y}_t = \rho I_{t-d}, \qquad d \in \{0,1,2,3\},
\]
where the delay $d$ is selected by validation MAE only. This model isolates the
contribution of observation delay without changing the latent epidemic process.

\subsection{Constrained structure discovery}

The discovery module searches over a bounded grammar rather than over arbitrary
symbolic dynamics. A structure specification contains:
\begin{itemize}[leftmargin=1.5em]
\item a compartment template from
\{\texttt{SIR}, \texttt{SEIR}, \texttt{SEIRS}, \texttt{SEIHR}, \texttt{SEIAR}\},
\item a fractional on/off flag,
\item an observation map from
\{\texttt{I}, \texttt{H}, \texttt{I+H}, \texttt{delayed\_I}\}, and
\item a delay parameter in $\{0,1,2,3\}$ when the observation map is
\texttt{delayed\_I}.
\end{itemize}

The search begins from \texttt{SEIR} and iterates a
propose-fit-verify-refine loop:
\begin{enumerate}[leftmargin=1.5em]
\item fit the current candidate;
\item generate one-hop neighbors via structure edits, observation-map edits,
fractional toggles, and delay edits;
\item reject invalid candidates using structural validity rules;
\item fit all remaining candidates with multiple random restarts;
\item score candidates using validation and rolling-origin criteria plus
complexity penalties; and
\item retain the top beam until convergence or early stopping.
\end{enumerate}

Validity rules enforce that infection originates from $S$, infectious-like
states eventually reach recovery, no isolated compartments exist, and
observation semantics are compatible with the latent template. For example,
\texttt{obs=H} and \texttt{obs=I+H} require an $H$ compartment, while
\texttt{obs=delayed\_I} requires an $I$ compartment and a valid delay.

To keep search tractable, delayed-observation neighbors are only introduced
from \texttt{SIR}, \texttt{SEIR}, and \texttt{SEIRS}; observation-map
neighbors for \texttt{SEIHR} prioritize \texttt{I}, \texttt{H}, and
\texttt{I+H}; and the maximum number of compartments remains capped at five.

\subsection{Discovery scoring}

Discovery is selected by a stability-aware score that combines predictive
quality and structural parsimony. At a high level, candidates are penalized for
\begin{itemize}[leftmargin=1.5em]
\item validation MAE,
\item rolling-origin error and stability,
\item additional free parameters,
\item additional compartments,
\item observation-map complexity,
\item delayed observation length, and
\item optional age-specific structural priors.
\end{itemize}

This design intentionally favors structures that are not only accurate on one
validation segment but also stable across rolling-origin forecasts. The age
prior can be toggled on or off as an explicit ablation, and the current
multi-seed observation-aware reruns show identical results with and without it.

\subsection{Benchmark-level uncertainty calibration}

Raw predictive intervals are produced only by the probabilistic SEIR model.
They are then calibrated \emph{after} benchmarking via a benchmark-level
post-processing pipeline. Five calibration kinds are compared:
\texttt{raw}, \texttt{scale\_calibrated}, \texttt{conformal\_absolute},
\texttt{conformal\_standardized}, and \texttt{conformal\_asymmetric}. Method
selection is based strictly on validation, never on test labels.

The current default uncertainty postprocessor is the balanced conformal
winner-selection rule \texttt{v3}, which trades a very small increase in
absolute coverage gap relative to the most coverage-oriented rule for visibly
better interval score and narrower widths.

\subsection{Robustness and ablations}

The current non-LLM evidence base includes:
\begin{itemize}[leftmargin=1.5em]
\item five-seed robustness analysis over seeds
$\{0,1,2,42,2024\}$,
\item a single-seed age-prior ablation,
\item an observation-aware five-seed no-age-prior ablation, and
\item conformal winner-rule comparison.
\end{itemize}

Together these controls allow us to separate three questions:
\begin{enumerate}[leftmargin=1.5em]
\item whether discovery wins are stable across optimization randomness;
\item whether delayed observation arises from data and the search objective
rather than from the age prior; and
\item whether probabilistic intervals can be calibrated without contaminating
point-forecast evaluation.
\end{enumerate}
